\section{Introduction}
Prime editing enables programmable substitutions, insertions and deletions without requiring a donor DNA template or a double-strand break \citep{anzalone2019}. Its practical success depends in part on the choice of prime-editing guide RNA (pegRNA). Alternative spacers, primer-binding sites (PBSs), reverse-transcription templates (RTTs) and guide constructs can produce different outcomes for the same intended edit. A useful computational method must therefore order candidate designs accurately and, when used to nominate a single candidate, favor designs with higher measured editing efficiency.

Existing methods have substantially advanced computational guide design. PRIDICT models editing efficiency and product purity; DeepPrime covers diverse editing systems and cellular settings; PRIDICT2.0 and ePRIDICT incorporate chromatin context; and OPED applies transfer learning to guide prediction and design \citep{mathis2023,yu2023,mathis2025,liu2023oped}. OptiPrime combines learned sequence-dependent quantities with a mechanistic description of prime-editing outcomes \citep{hsu2026}. These approaches provide a strong foundation for testing whether an edit-aware representation and order-sensitive supervision can further improve pegRNA ranking. Here, released OptiPrime is the primary comparator on a shared source benchmark and an external library; our comparisons do not reevaluate all applications reported in its original study.

We propose \model{}, a sequence-modeling framework designed to integrate three sources of information relevant to candidate ordering: the intended nucleotide change, the structure of the pegRNA and the experimental context. One encoder processes aligned wild-type/edited base pairs, while a second processes the spacer, PBS and RTT with segment identity. Bidirectional cross-attention connects the two streams, and context conditioning accounts for cell, editor and construct metadata. We investigate self-attention and diagonal state-space mixing within this framework. Its ordinal head learns whether a measured outcome exceeds a series of training-derived thresholds, producing an order-sensitive score without requiring direct estimation of a precisely calibrated efficiency. This supervision is complemented by proportion-based heads in the source ensemble. The contribution is their integration and empirical evaluation for pegRNA ranking, rather than a claim that these individual neural-network components are new.

We assess \model{} in two stages. First, we test rank-based performance using pooled Spearman correlation, within-condition and within-protospacer correlations, Kendall's tau-b and external-library transfer. These complementary analyses distinguish broad ordering accuracy from gains driven only by differences between experimental conditions or by zero-valued outcomes. Second, we evaluate the measured efficiency of the model's highest-ranked candidate within a group that fixes the intended allele and experimental context. This selection endpoint directly reflects the design decision, whereas even a strong pooled correlation need not guarantee the best ordering within every candidate set. Absolute-efficiency calibration is treated as a secondary property, not the paper's primary objective.

Our results show that \model{} outperforms OptiPrime across the reported source rank-based metrics and retains a pooled ranking advantage on the external library. It also improves fixed-edit selection on the source benchmark. On the external library, native-ordinal fine-tuning yields higher top-choice efficiency than released OptiPrime across three optimizer seeds, although the zero-shot model does not. We therefore distinguish source-trained ranking performance from target-supervised selection performance throughout. Additional analyses examine architecture and training choices, alternative adaptation heads, label budgets and source retention. All research is computational, using previously measured outcomes; later adaptation experiments reuse the external library and are not presented as independent-study confirmation.
